\clearpage
\begin{examplebox}
\textbf{\textcolor{CExample}{先做一道“发现规则”的例题：为什么要解括号内的不等式？}}

\textbf{题目：} $f(t)=\sqrt{t-1}+\sqrt{3-t}$。
先求 $f(t)$ 的定义域，再求 $f(2x-1)$ 的定义域，最后尝试不写具体的 $f$ 来求。

\textbf{第一问（根号原来的要求）：} $t-1\ge0$、$3-t\ge0$，
于是 $f(t)$ 允许 $t\in[1,3]$。

\textbf{第二问（把数真的放进去）：}
\[
 f(2x-1)=\sqrt{2x-2}+\sqrt{4-2x}.
\]
它有意义时，$2x-2\ge0$、$4-2x\ge0$，所以 $x\in[1,2]$。

\textbf{第三问（把具体式子遮住）：} 现在只知道“$f$ 接受 $[1,3]$ 的数”，
但 $f(2x-1)$ 要把 $2x-1$ 送进去，便必须有
\[
 1\le2x-1\le3
 \quad\Longleftrightarrow\quad \boxed{1\le x\le2}.
\]
\textbf{发现：} 结果和第二问完全一样！
\textbf{具体计算产生了规则；知道输入要求后，就不必再展开具体式子。}

\medskip

\textbf{\textcolor{CExample}{例1：已知 $f(x)$ 的定义域，直接求新定义域（基础）}}

\textbf{题目：} 已知 $f(t)$ 的定义域是 $[0,3]$，求 $f(2x+1)$ 的定义域。

\textbf{第1步：} $f$ 只接受 $0$ 到 $3$，所以括号里的 $2x+1$ 必须在此范围内：
\[
 0\le2x+1\le3.
\]
\textbf{第2步：} 各部分同时减去 $1$：$-1\le2x\le2$。
再同时除以 $2$，得到
\[
 -\frac12\le x\le1.
\]
\textbf{答案：} $\boxed{[-\frac12,1]}$。

\textbf{核对：} 取左端点 $x=-\frac12$，括号内是 $0$；
取右端点 $x=1$，括号内是 $3$，都允许。

\medskip
\textbf{\textcolor{CExample}{例2：已知 $f(2x-1)$ 的定义域，反求 $f(t)$（基础）}}

\textbf{题目：} $f(2x-1)$ 的完整定义域是 $[1,4]$，求 $f(t)$ 的定义域。

\textbf{第1步：} 题目给的是\textbf{$x$ 的范围} $[1,4]$，不是 $2x-1$ 的范围。

\textbf{第2步：} 当 $x$ 从 $1$ 到 $4$ 时，
\[
 1\le2x-1\le7.
\]
而且 $2x-1$ 在两端之间的每个数都能取到。

\textbf{第3步：} 因为 $2x-1$ 随着 $x$ 的变化能取到任意实数，
题目给出的完整定义域足以确定 $f$ 全部允许的输入数。

\textbf{答案：} $\boxed{[1,7]}$。

\medskip
\textbf{\textcolor{CExample}{例3：先找 $f$ 的输入范围，再求另一个式子（常见综合题）}}

\textbf{题目：} 已知 $f(x+1)$ 的完整定义域是 $[0,2]$，求 $f(2x-1)$ 的定义域。

\textbf{第1步：} 由 $0\le x\le2$，两边加 $1$ 得
\[
 1\le x+1\le3.
\]
因为 $x+1$ 能取到任意实数，所以 $f$ 的定义域恰好是 $[1,3]$。

\textbf{第2步：} 现在求 $f(2x-1)$，只需让 $2x-1$ 落入 $[1,3]$：
\[
 1\le2x-1\le3.
\]
\textbf{第3步：} 加 $1$ 再除以 $2$：$1\le x\le2$。

\textbf{答案：} $\boxed{[1,2]}$。

\textbf{速记：} $f(x+1)$ 的 $x$ 范围 $\to$ $f$ 真正接受的数 $\to$ $f(2x-1)$ 的 $x$ 范围。

\medskip
\begin{minipage}{\linewidth}
\textbf{\textcolor{CExample}{例4：有负号、有开端点，怎么办？（易考）}}
\begin{center}\begin{tikzpicture}[font=\small,>=stealth,line join=round]
\draw[gray!45,rounded corners=3pt] (0,0) rectangle (14.2,9.4);
\node at (7.1,8.8) {负系数使端点次序反转，开闭性质仍跟随原点};

\begin{scope}[shift={(3.2,5.8)},x=1.25cm,y=.6cm]
\draw[->] (-.5,0)--(3.7,0) node[right]{$x$};\draw[->] (0,-5.9)--(0,2) node[above]{$t$};
\draw[teal!80!black,thick,domain=-.2:3.3] plot(\x,{1-2*\x});
\draw[dashed] (3,0)--(3,-5)--(0,-5);\draw[orange!90!black,line width=2pt] (0,0)--(3,0);\draw[blue,line width=2pt] (0,-5)--(0,1);
\fill (3,-5) circle(1.6pt);\draw[thick,fill=white] (0,1) circle(2pt);\fill (3,0) circle(1.6pt);\draw[thick,fill=white] (0,0) circle(2pt);
\node[above] at (3,0){$3$};\node[left] at (0,-5){$-5$};\node[left] at (0,1){$1$};
\end{scope}
\node[align=left,anchor=west] at (8.2,5.1){$t=1-2x$\\[10pt] $x\in(0,3]$\\[8pt] $t\in[-5,1)$\\[10pt] 空心：不取到\\[4pt] 实心：取到};
\node[text width=13.2cm,align=center] at (7.1,.65) {改变变量时，数值大小顺序会变，端点是否取到不会丢失。};
\end{tikzpicture}
\end{center}
\end{minipage}



\textbf{题目：} $f(1-2x)$ 的完整定义域为 $(0,3]$，求 $f(x+1)$ 的定义域。

\textbf{第1步：} 已知 $0<x\le3$。乘以 $-2$，不等号方向改变：
\[
 -6\le -2x<0.
\]
再加 $1$，得到
\[
 -5\le1-2x<1.
\]
所以 $f(t)$ 的定义域是 $[-5,1)$。

\textbf{第2步：} 让新括号里的 $x+1$ 落入 $[-5,1)$：
\[
 -5\le x+1<1.
\]
两边减 $1$，得到 $-6\le x<0$。

\textbf{答案：} $\boxed{[-6,0)}$。

\textbf{核对：} $x=0$ 不能取，因为此时新括号内是 $1$，而 $f$ 不接受 $1$。

\medskip
\textbf{\textcolor{CExample}{例5：已知 $f(x^2)$ 的定义域，能反求 $f(x)$ 吗？（眼前一亮）}}

\textbf{题目：} 已知 $f(x^2)$ 的完整定义域是 $[-2,2]$，问 $f(t)$ 的定义域是不是 $[0,4]$？

\textbf{第1步：} 在 $[-2,2]$ 内，$x^2$ 恰好能取遍 $[0,4]$，
所以 $f$ 必须允许 $[0,4]$ 中的每一个数。

\textbf{第2步：} 因为 $|x|>2$ 时复合函数没有意义，
所有大于 $4$ 的输入数都不被 $f$ 允许。

\textbf{第3步：} 但是 $x^2$ 根本送不出负数，
题目无法告诉我们 $f$ 是否还接受负数。

\textbf{答案：} \textbf{不能唯一确定} $f(t)$ 的完整定义域。
只能确定它在非负数范围内是 $[0,4]$。
例如 $[0,4]$ 和 $(-\infty,4]$ 都与题目相容。

\medskip
\textbf{\textcolor{CExample}{例6：括号里的式子本身也可能有限制（补充）}}

\textbf{题目：} 已知 $f(t)$ 的定义域是 $[-2,2]$，求 $f(\sqrt{x-1})$ 的定义域。

\textbf{第1步：} 先保证根号有意义，$x-1\ge0$，所以 $x\ge1$。

\textbf{第2步：} 根号的值永远不小于 $0$，因此只需它不超过 $2$：
\[
 0\le\sqrt{x-1}\le2.
\]
\textbf{第3步：} 两边平方（这里两边都不小于 $0$），得到
$x-1\le4$，即 $x\le5$。
合并第1步的条件，得
\[
 \boxed{1\le x\le5}.
\]
\textbf{提醒：} 若忘了 $x\ge1$，即使后面计算正确，也不能算完整答案。
\end{examplebox}
